The principal results in Chapter 1 summarize the calculation. They are
not assumptions in the derivation that starts here. This chapter derives
the Compton amplitudes and their parameter integrals
from the QED vertices. The loop, energy, and angular integrations are performed on these
amplitudes. The total in Equation (6) of Lee, Schwartz and
Zhang~\cite{qed-processes-lsz} is reserved for a final comparison. No coefficient from that total is used
to determine an integral or its integration constant here.

\section{Conventions and diagram counting}
Let $p,k$ denote incoming electron and photon momenta, and $p',k'$
the outgoing momenta. The electron mass is $m$, its charge is
$-e_{\mathrm T}$, and $\alpha_{\mathrm T}=e_{\mathrm T}^2/(4\pi)$
is the Thomson coupling. Use the mostly-minus metric, the vertex
$+\ii e_{\mathrm T}\gamma^\mu$, and
\begin{align}
 p+k&=p'+k',\qquad p^2=p'^2=m^2,\qquad k^2=k'^2=0,\\
 s&=(p+k)^2,\qquad t=(p-p')^2,\qquad u=(p-k')^2.
\end{align}
We write $\epsilon=\epsilon_{\rm UV}$ and regulate ultraviolet integrals in $d=4-2\epsilon$ dimensions,
with coupling $\mu^\epsilon e_{\mathrm T}$, and soft photons by a
mass $\lambda$. The propagators are
\begin{equation}
 S_F(P)=\ii S_0(P),\qquad S_0(P)=\frac{\slashed P+m}{P^2-m^2+\ii0},
 \qquad D_{\rho\sigma}(\ell)=\frac{-\ii\eta_{\rho\sigma}}
                                      {\ell^2-\lambda^2+\ii0}.
\end{equation}
The photon symbol omits $F$ but has Feynman boundary conditions.
External photons remain massless.
The Born amplitude is
\begin{equation}
 \mathcal M_0=-e_{\mathrm T}^2\epsilon_\mu\epsilon_\nu'^*
 \bar u(p')R_0^{\nu\mu}u(p),\qquad
 R_0^{\nu\mu}=\gamma^\nu S_0(p+k)\gamma^\mu
                         +\gamma^\mu S_0(p-k')\gamma^\nu.
\end{equation}
At one loop the electron line has four photon vertices. Choosing
the positions of the two external photons gives six choices and
ordering those photons gives two choices. Of these twelve diagrams,
four are external-electron self-energy diagrams. On-shell external
normalization accounts for those four. The remaining eight are two
internal self-energies, four vertices and two spanning photons.
There is no exchanged Born photon on which to insert a vacuum
polarization. The physical charge and real-photon normalization
account for the photon two-point function at zero momentum squared.

\section{Tensor loop integration}
Define an integral whose factor $1/\ii$ removes the Minkowski loop
factor. Here $\Gamma$ denotes Euler's gamma function.
\begin{equation}
 I_n(\Delta)=\frac{\mu^{2\epsilon}}{\ii}
 \int\frac{\dd^dL}{(2\pi)^d}\frac1{(L^2-\Delta+\ii0)^n}
 =\frac{(-1)^n\mu^{2\epsilon}}{(4\pi)^{d/2}}
 \frac{\Gamma(n-d/2)}{\Gamma(n)}(\Delta-\ii0)^{d/2-n}.
 \label{eq:comp-supp-master}
\end{equation}
To derive the last expression, first take $\Delta>0$ and
$\operatorname{Re}n>\operatorname{Re}d/2$. A Wick rotation supplies
$\ii(-1)^n$, and the Euclidean angular integral followed by
$v=L_E^2$ gives
\begin{equation}
 I_n(\Delta)=\frac{(-1)^n\mu^{2\epsilon}}
 {(4\pi)^{d/2}\Gamma(d/2)}
 \int_0^\infty\frac{v^{d/2-1}\dd v}{(v+\Delta)^n}.
\end{equation}
Scaling $v=\Delta z$ reduces the remaining integral to Euler's identity
\begin{equation}
 \int_0^\infty\frac{z^{a-1}\dd z}{(1+z)^{a+b}}
 =\frac{\Gamma(a)\Gamma(b)}{\Gamma(a+b)},
 \qquad \operatorname{Re}a,\operatorname{Re}b>0,
\end{equation}
with $a=d/2$ and $b=n-d/2$. Analytic continuation in $d$ and in
$\Delta-\ii0$ then gives \eqref{eq:comp-supp-master} with its causal
boundary value. Odd powers of $L$ vanish. Rotational invariance of the
regulated integral gives
\begin{align}
 \int\dd^dL\,L_\rho L_\sigma f(L^2)
 &=\frac{\eta_{\rho\sigma}}d\int\dd^dL\,L^2f(L^2),\\
 \frac{L^2}{(L^2-\Delta)^n}
 &=\frac1{(L^2-\Delta)^{n-1}}
                  +\frac{\Delta}{(L^2-\Delta)^n}.
\end{align}
The scalar integral therefore also determines the quadratic
numerators. Combine $n$ denominators $D_j$ with
\begin{equation}
 \frac1{\prod_{j=1}^nD_j}
 =\Gamma(n)\int_{a_j\geq0}\prod_j\dd a_j\,
 \frac{\delta(1-\sum_j a_j)}{(\sum_j a_jD_j)^n}.
 \label{eq:comp-supp-feynman-parameters}
\end{equation}
The simplex measure is normalized to $1/(n-1)!$.
This normalization is needed when checking pole coefficients.

\section{Self-energy and residue subtraction}
The insertion $-\ii\Sigma(P)$ yields
\begin{equation}
 \Sigma(P)=\frac{e_{\mathrm T}^2\mu^{2\epsilon}}{\ii}
 \int\frac{\dd^d\ell}{(2\pi)^d}
 \frac{\gamma^\rho(\slashed P-\slashed\ell+m)\gamma_\rho}
 {[(P-\ell)^2-m^2+\ii0](\ell^2-\lambda^2+\ii0)}.
\end{equation}
Use $v$ for the electron denominator. Completing the square gives
\begin{equation}
 v[(P-\ell)^2-m^2]+(1-v)(\ell^2-\lambda^2)
 = (\ell-vP)^2-\Delta_\Sigma,
 \quad \Delta_\Sigma=vm^2+(1-v)\lambda^2-v(1-v)P^2.
\end{equation}
After the shift $L=\ell-vP$, the gamma identities
$\gamma^\rho\gamma^\alpha\gamma_\rho=(2-d)\gamma^\alpha$
and $\gamma^\rho\gamma_\rho=d$ leave
\begin{equation}
 \Sigma(P)=\frac{e_{\mathrm T}^2\mu^{2\epsilon}}{(4\pi)^{d/2}}
 \Gamma(\epsilon)\int_0^1\dd v\,
 \frac{(2-d)(1-v)\slashed P+dm}{(\Delta_\Sigma-\ii0)^\epsilon}
 =a(P^2)\slashed P+mb(P^2).
\end{equation}
For an explicit expansion define
\begin{equation}
 C=\frac{\alpha_{\mathrm T}}{4\pi},\qquad
 \mathcal L(\Delta)=\frac1\epsilon-\gamma_E+\ln4\pi
                         +\ln\frac{\mu^2}{\Delta-\ii0},
\end{equation}
where $\gamma_E$ is Euler's constant. Through the finite terms,
\begin{align}
 a(P^2)&=C\int_0^1\dd v\,
       [-2(1-v)\mathcal L(\Delta_\Sigma)+2(1-v)],\\
 b(P^2)&=C\int_0^1\dd v\,[4\mathcal L(\Delta_\Sigma)-2].
\end{align}
The last terms come from multiplying $d-4$ by the ultraviolet pole.
They would be lost by doing the numerator algebra in four dimensions
before subtraction.

To determine the pole and its residue, regard $\slashed P$ as a
variable $w$ near $m$. Then
\begin{equation}
 \begin{aligned}
 a(w^2)w+mb(w^2)={}&m[a(m^2)+b(m^2)]\\
 &+(w-m)\{a(m^2)+2m^2[a'(m^2)+b'(m^2)]\}\\
 &+\Order((w-m)^2).
 \end{aligned}
\end{equation}
Here the primes differentiate the scalar squared argument. Subtract
the first two coefficients,
\begin{align}
 c_m&=m[a(m^2)+b(m^2)],\qquad
 c_Z=a(m^2)+2m^2[a'(m^2)+b'(m^2)],\\
 \Sigma_{\mathrm R}(P)&=\Sigma(P)-c_m-(\slashed P-m)c_Z.
\end{align}
The renormalized propagator has a pole at $m$ with unit residue.
Writing $\Delta_0=v^2m^2+(1-v)\lambda^2$ gives explicit subtraction
integrals,
\begin{align}
 c_m&=mC\int_0^1\dd v\,[2(1+v)\mathcal L(\Delta_0)-2v],\\
 c_Z&=C\int_0^1\dd v\,
 \left[-2(1-v)\mathcal L(\Delta_0)+2(1-v)
       +\frac{4m^2v(1-v)(1+v)}{\Delta_0}\right].
 \label{eq:comp-supp-cz}
\end{align}
At fixed $\lambda>0$ these are ordinary one-dimensional parameter
integrals. The last term in $c_Z$ displays the logarithmic soft
sensitivity as $\lambda\to0$.

\subsection{Evaluation of the on-shell constants}
The logarithmic integrals in the mass subtraction are integrable at
$v=0$, whereas the rational term in the residue subtraction is not
integrable there after setting the photon mass to zero. To separate
these two cases, put
\begin{equation}
 r_\lambda=\frac{\lambda}{m},\qquad
 d_0(v)=v^2+(1-v)r_\lambda^2,\qquad
 U_m=\frac1\epsilon-\gamma_E+\ln4\pi+\ln\frac{\mu^2}{m^2}.
\end{equation}
Then $\mathcal L(\Delta_0)=U_m-\ln d_0(v)$. For the logarithmic
terms the limit can be taken under the integral, giving
\begin{align}
 \frac{c_m}{mC}
 &=3U_m-4\int_0^1(1+v)\ln v\,\dd v-1+o(1)\\
 &=3U_m+4+o(1).
\end{align}
For the regular part of $c_Z/C$, define and evaluate
\begin{align}
 I_{\rm reg}&=\int_0^1[-2(1-v)\mathcal L(\Delta_0)+2(1-v)]\,\dd v\\
 &=-U_m+4\int_0^1(1-v)\ln v\,\dd v+1+o(1)\\
 &=-U_m-2+o(1).
\end{align}
Here $o(1)$ means a term vanishing as $r_\lambda\to0$ at fixed
$\mu/m$. The elementary moments used are
$\int_0^1\ln v\,\dd v=-1$ and
$\int_0^1v\ln v\,\dd v=-1/4$.

The remaining term must be integrated before this limit. Replacing
$d_0(v)$ by $v^2+r_\lambda^2$ changes its integral by $O(r_\lambda)$.
For example, split the interval at $v=r_\lambda$ and $v=1/2$.
The denominator difference is $v r_\lambda^2$, so the absolute
integrand difference is bounded by a constant times
$r_\lambda^2v^2/(v^2+r_\lambda^2)^2$ on the first two intervals
and by a constant times $r_\lambda^2$ on the last. Its integral
therefore vanishes. Consequently,
\begin{align}
 4\int_0^1\frac{v(1-v^2)}{d_0(v)}\,\dd v
 &=4\int_0^1\frac{v-v^3}{v^2+r_\lambda^2}\,\dd v+o(1)\\
 &=2(1+r_\lambda^2)\ln\frac{1+r_\lambda^2}{r_\lambda^2}-2+o(1)\\
 &=-4\ln r_\lambda-2+o(1).
\end{align}
Combining the regular and singular parts fixes the finite constant
as well as the infrared logarithm,
\begin{equation}
 c_Z=-C[U_m+4+4\ln r_\lambda]+C\,o(1).
\end{equation}

To relate these subtractions to the Lagrangian, write
$Z_j=1+\delta_j+O(\alpha_{\mathrm T}^2)$,
$\psi_0=Z_2^{1/2}\psi$, and $m_0=Z_m m$. The electron quadratic
counterterm is
\begin{equation}
 \delta\mathcal L_e=
 \overline\psi[\delta_2(\ii\slashed\partial-m)-m\delta_m]\psi.
\end{equation}
Its insertion is $\ii[\delta_2(\slashed P-m)-m\delta_m]$.
Combining it with $-\ii\Sigma$ gives
$\Sigma_R=\Sigma-\delta_2(\slashed P-m)+m\delta_m$.
Comparison with the pole and residue subtractions thus determines
\begin{equation}
 \delta_m=-\frac{c_m}{m},\qquad \delta_2=c_Z.
\end{equation}
This also explains why the field counterterm multiplies both the
kinetic term and the renormalized mass term.

These mass and field counterterms agree with the diagonal QED limit
of Denner's on-shell expressions, Equation (3.20) of
Reference~\cite{compton-overview-denner-renormalization}.
In his self-energy convention the left- and right-chiral scalar
coefficients both equal $-a$, and the mass coefficient equals $-b$.
His additive mass counterterm therefore equals $-m(a+b)=m\delta_m$,
and either diagonal field counterterm equals
$a+2m^2(a'+b')=\delta_2$. This comparison is made after evaluating
our parameter integrals.

\section{The off-shell vertex}
At a photon insertion let the electron momenta be $P$ and $P'$,
with $q=P'-P$. The correction to the dimensionless vertex is
\begin{equation}
 V^\rho(P',P)=\frac{e_{\mathrm T}^2\mu^{2\epsilon}}{\ii}
 \int\frac{\dd^d\ell}{(2\pi)^d}
 \frac{\gamma^\alpha S_0(P'-\ell)\gamma^\rho
                         S_0(P-\ell)\gamma_\alpha}
      {\ell^2-\lambda^2+\ii0}.
\end{equation}
Introduce $x,y,z\geq0$, $x+y+z=1$, for the two electron lines and
the photon. With $\int_{\mathcal S_2}=\int_0^1\dd x\int_0^{1-x}\dd y$,
\begin{equation}
 K=xP'+yP,\qquad
 \Delta=(x+y)m^2+z\lambda^2-xzP'^2-yzP^2-xyq^2.
\end{equation}
Let $A=\slashed P'-\slashed K+m$ and $B=\slashed P-\slashed K+m$.
After shifting $\ell=L+K$, the numerator is
\begin{equation}
 N^\rho(L+K)=\gamma^\alpha(A-\slashed L)\gamma^\rho
 (B-\slashed L)\gamma_\alpha
 \longrightarrow N^\rho(K)+\frac{(d-2)^2}{d}L^2\gamma^\rho.
\end{equation}
The arrow denotes symmetric loop integration. Equations
\eqref{eq:comp-supp-master} and \eqref{eq:comp-supp-feynman-parameters}
then give
\begin{equation}
 V^\rho=\frac{e_{\mathrm T}^2\mu^{2\epsilon}}{(4\pi)^{d/2}}
 \int_{\mathcal S_2}
 \left[-\frac{\Gamma(1+\epsilon)N^\rho(K)}{(\Delta-\ii0)^{1+\epsilon}}
 +\frac{(d-2)^2\Gamma(\epsilon)\gamma^\rho}
 {2(\Delta-\ii0)^\epsilon}\right].
\end{equation}
Equivalently, its finite expansion is
\begin{equation}
 V^\rho=C\int_{\mathcal S_2}
 \left[-\frac{N^\rho_4(K)}{\Delta-\ii0}
             +\gamma^\rho\{2\mathcal L(\Delta)-4\}\right].
 \label{eq:comp-supp-vertex-finite}
\end{equation}
The subscript $4$ permits four-dimensional gamma algebra in the
nondivergent numerator term. Add the counterterm
$V_{\mathrm R}^\rho=V^\rho+c_Z\gamma^\rho$.
The pole of $V$ is $C\gamma^\rho/\epsilon$, since the simplex has
area $1/2$. The pole of $c_Z$ is $-C/\epsilon$. Their cancellation
also removes the common ultraviolet scale dependence.

The relative sign follows independently from
\begin{equation}
 S_0(A)(\slashed A-\slashed B)S_0(B)=S_0(B)-S_0(A),
 \qquad q_\rho V^\rho(P',P)=\Sigma(P)-\Sigma(P').
\end{equation}
On adding the counterterms,
\begin{equation}
 q_\rho V_{\mathrm R}^\rho(P',P)
 =\Sigma_{\mathrm R}(P)-\Sigma_{\mathrm R}(P').
 \label{eq:comp-supp-ward}
\end{equation}
A decomposition into only the on-shell Dirac and Pauli form factors
cannot satisfy this off-shell identity for general $P,P'$.

\subsection{Photon normalization and the charge counterterm}
The off-shell Ward identity fixes the vertex counterterm to the same
$c_Z$ as the electron field counterterm. In the notation
$Z_j=1+\delta_j+O(\alpha_{\mathrm T}^2)$ this gives
\begin{equation}
 \delta_1=\delta_2=c_Z.
\end{equation}
To determine the charge counterterm, define
$A_{0\mu}=Z_3^{1/2}A_\mu$ and
$e_0=\mu^\epsilon Z_e e_{\mathrm T}$. Substituting these and the
bare electron field into the interaction gives
\begin{equation}
 Z_1=Z_e Z_2 Z_3^{1/2},\qquad
 \delta_1=\delta_e+\delta_2+\tfrac12\delta_3,
 \qquad \delta_e=-\tfrac12\delta_3.
\end{equation}
Thus only the photon residue remains to be calculated.

For one electron species, the closed electron loop contributes
\begin{equation}
 \ii\Pi^{\mu\nu}(q)=-e_{\mathrm T}^2\mu^{2\epsilon}
 \int\frac{\dd^d\ell}{(2\pi)^d}
 \frac{\operatorname{tr}[\gamma^\mu(\slashed\ell+m)
           \gamma^\nu(\slashed\ell+\slashed q+m)]}
 { (\ell^2-m^2+\ii0)[(\ell+q)^2-m^2+\ii0]}.
\end{equation}
The sign includes the minus sign of a closed fermion loop. The trace is
\begin{equation}
 4\{\ell^\mu(\ell+q)^\nu+\ell^\nu(\ell+q)^\mu
       -\eta^{\mu\nu}[\ell\cdot(\ell+q)-m^2]\}.
\end{equation}
Combine the denominators with a parameter $v$ and shift
$L=\ell+vq$. Their quadratic form is
$L^2-\Delta_\gamma+\ii0$, where
$\Delta_\gamma=m^2-v(1-v)q^2$.
Odd terms vanish, and symmetric integration replaces
$L^\mu L^\nu$ by $\eta^{\mu\nu}L^2/d$. The remaining terms
proportional to the metric are reduced using
\begin{equation}
 0=\int\dd^dL\,\frac{\partial}{\partial L^\rho}
       \frac{L^\rho}{L^2-\Delta_\gamma+\ii0}
 =\int\dd^dL\left[\frac{d}{L^2-\Delta_\gamma+\ii0}
       -\frac{2L^2}{(L^2-\Delta_\gamma+\ii0)^2}\right].
\end{equation}
This identity follows in a convergent domain of dimensional
regularization and extends by analytic continuation. Together with
the scalar loop integral it gives the transverse form
\begin{align}
 \ii\Pi^{\mu\nu}(q)
 &=\ii(q^\mu q^\nu-q^2\eta^{\mu\nu})\Pi_\gamma(q^2),\\
 \Pi_\gamma(q^2)
 &=8C\int_0^1v(1-v)\mathcal L(\Delta_\gamma)\,\dd v
       +O(\epsilon),\\
 \Pi_\gamma(0)&=\frac{4C}{3}U_m.
\end{align}
There is no soft photon in this loop, so its zero-momentum value
has no photon-mass logarithm. With this sign convention the transverse
quadratic action has coefficient $1+\delta_3+\Pi_\gamma(q^2)$.
Unit residue at $q^2=0$ therefore requires
\begin{equation}
 \delta_3=-\Pi_\gamma(0)=-\frac{\alpha_{\mathrm T}}{3\pi}U_m,
 \qquad
 \delta_e=\frac{\alpha_{\mathrm T}}{6\pi}U_m.
\end{equation}
These relations agree with the QED limit of Denner's photon and
charge counterterms, Equations (3.19) and (3.32) of
Reference~\cite{compton-overview-denner-renormalization}.
For a Compton amplitude the two charge factors and the two external
photon field factors combine to $Z_e^2Z_3=1$ at this order.
This cancellation accounts for photon normalization even though
there is no internal photon propagator in the Born amplitude.

\section{The spanning photon integral}
Put $Q=p+k$. The remaining $s$-ordered matrix is
\begin{equation}
 W_s^{\nu\mu}=\frac{e_{\mathrm T}^2}{\ii}
 \int\frac{\dd^4\ell}{(2\pi)^4}
 \frac{N^{\nu\mu}(\ell)}
 {D_0D_1D_2D_3},
\end{equation}
where
\begin{align}
 (D_0,D_1,D_2,D_3)&=(\ell^2-\lambda^2,
 (p-\ell)^2-m^2,(Q-\ell)^2-m^2,(p'-\ell)^2-m^2)+\ii0,\\
 N^{\nu\mu}(\ell)&=\gamma^\alpha(\slashed p'-\slashed\ell+m)
 \gamma^\nu(\slashed Q-\slashed\ell+m)
 \gamma^\mu(\slashed p-\slashed\ell+m)\gamma_\alpha.
\end{align}
Each component of the denominator tuple carries $+\ii0$.
The loop is ultraviolet convergent by power counting. For simplex
parameters $a_0,\ldots,a_3$ define
\begin{equation}
 K=a_1p+a_2Q+a_3p',\qquad
 \Delta=K^2+a_0\lambda^2+a_2(m^2-s),\qquad
 \sum_j a_jD_j=(\ell-K)^2-\Delta+\ii0.
\end{equation}
Expand the cubic numerator about $K$. Only its constant and second
derivative contribute,
\begin{equation}
 N(L+K)\longrightarrow N(K)+\frac{L^2}{8}\Box_\ell N(K),
 \qquad \Box_\ell=\frac{\partial}{\partial\ell^\rho}
                       \frac{\partial}{\partial\ell_\rho}.
\end{equation}
The factor $1/8$ is the Taylor factor $1/2$ times the tensor
average $1/4$. More explicitly, with $A_1=\slashed p-\slashed K+m$,
$A_2=\slashed Q-\slashed K+m$, $A_3=\slashed p'-\slashed K+m$,
\begin{equation}
 \begin{aligned}
 \Box_\ell N(K)=2\gamma^\alpha\big[&
 \gamma^\rho\gamma^\nu\gamma_\rho\gamma^\mu A_1
 +\gamma^\rho\gamma^\nu A_2\gamma^\mu\gamma_\rho\\
 &+A_3\gamma^\nu\gamma^\rho\gamma^\mu\gamma_\rho
 \big]\gamma_\alpha.
 \end{aligned}
\end{equation}
There are three ways to differentiate two different electron
numerators. All have a positive sign because each derivative brings
a minus sign. The needed scalar integrals at $d=4$ are
\begin{equation}
 3! I_4=\frac1{16\pi^2\Delta^2},\qquad
 3!(I_3+\Delta I_4)=\frac{-2}{16\pi^2\Delta}.
\end{equation}
Thus the complete parameter representation is
\begin{equation}
 W_s^{\nu\mu}=C\int_{\mathcal S_3}\dd^4a
 \left[\frac{N^{\nu\mu}(K)}{(\Delta-\ii0)^2}
       -\frac{\Box_\ell N^{\nu\mu}(K)}{4(\Delta-\ii0)}\right],
 \qquad \int_{\mathcal S_3}\dd^4a=
 \int_0^1\prod_{j=0}^3\dd a_j\,\delta(1-\sum_j a_j).
\end{equation}
Physical $s$-channel cuts require the displayed complex boundary
value. Replacing it by a positive real denominator would remove
physical absorptive terms.

\section{Assembly and external Ward identities}
The $s$ ordering is
\begin{equation}
 \begin{aligned}
 R_{1,s}^{\nu\mu}={}&V_{\mathrm R}^\nu(p',Q)S_0(Q)\gamma^\mu
 +\gamma^\nu S_0(Q)V_{\mathrm R}^\mu(Q,p)\\
 &+\gamma^\nu S_0(Q)\Sigma_{\mathrm R}(Q)S_0(Q)\gamma^\mu
 +W_s^{\nu\mu}.
 \end{aligned}
\end{equation}
Interchange $(\mu,k)$ and $(\nu,-k')$ to obtain $R_{1,u}$, and
set $R_1=R_{1,s}+R_{1,u}$. Its amplitude is
$\mathcal M_1=-e_{\mathrm T}^2\epsilon_\mu\epsilon_\nu'^*
\bar u(p')R_1^{\nu\mu}u(p)$.
For the incoming Ward identity, write $\doteq$ for equality between
external on-shell spinors. The four $s$-ordered terms contract to
\begin{align}
 k_\mu V_{\mathrm R}^\nu S_0(Q)\gamma^\mu
 &\doteq V_{\mathrm R}^\nu(p',Q),\\
 k_\mu\gamma^\nu S_0(Q)V_{\mathrm R}^\mu(Q,p)
 &\doteq-\gamma^\nu S_0(Q)\Sigma_{\mathrm R}(Q),\\
 k_\mu\gamma^\nu S_0(Q)\Sigma_{\mathrm R}(Q)S_0(Q)\gamma^\mu
 &\doteq\gamma^\nu S_0(Q)\Sigma_{\mathrm R}(Q),\\
 k_\mu W_s^{\nu\mu}&=V^\nu(p',p)-V^\nu(p',Q).
\end{align}
Their sum is $V_{\mathrm R}^\nu(p',p)$. The $u$ ordering yields
its negative. Hence $k_\mu R_1^{\nu\mu}\doteq0$. The outgoing
identity $k'_\nu R_1^{\nu\mu}\doteq0$ follows in the same way.
In particular, omitting the spanning graph spoils the cancellation.

Define $\overline A=\gamma^0A^\dagger\gamma^0$ and the bilinear
\begin{equation}
 \mathcal T(A,B)=\frac14\eta_{\mu\mu'}\eta_{\nu\nu'}
 \operatorname{tr}[(\slashed p'+m)A^{\nu\mu}
 (\slashed p+m)\overline B^{\nu'\mu'}].
 \label{eq:comp-supp-trace-bilinear}
\end{equation}
The factor $1/4$ averages over the two incoming electron spins and
photon polarizations. For amplitudes satisfying both external Ward
identities, the two physical polarization sums can be replaced by
$-\eta_{\mu\mu'}$ and $-\eta_{\nu\nu'}$. Their product gives the
positive metric contraction displayed here. The virtual interference
is $V=2e_{\mathrm T}^4\operatorname{Re}\mathcal T(R_1,R_0)$,
where the factor two comes from the two cross terms in the squared
amplitude. Explicitly,
\begin{align}
 V&=\frac{e_{\mathrm T}^4}{2}\operatorname{Re}
 \eta_{\mu\mu'}\eta_{\nu\nu'}\operatorname{tr}
 [(\slashed p'+m)R_1^{\nu\mu}(\slashed p+m)
                         \overline R_0^{\nu'\mu'}],\\
 B_{\mathrm{amp}}&=2e_{\mathrm T}^4
 \left(\frac{\omega}{\omega'}+\frac{\omega'}{\omega}-\sin^2\theta\right),
 \qquad \omega'=\frac{\omega}{1+(\omega/m)(1-\cos\theta)},\\
 \frac{\dd\sigma_{\mathrm{Born+virt}}}{\dd\Omega}
 &=\frac1{64\pi^2m^2}\left(\frac{\omega'}{\omega}\right)^2(B_{\mathrm{amp}}+V).
\end{align}
The rest-frame incident energy is $\omega=k^0$, and $\theta$
is the angle between the incoming and outgoing photons in that frame.
Here $B_{\mathrm{amp}}=\overline{|\mathcal M_0|^2}$ is the Born squared amplitude,
including its coupling factors.
These expressions average over both initial spin and initial photon
polarization. No square of $R_1$ belongs at this order.

\subsection{Integration of the Born cross section}
The Born total follows directly from the preceding differential cross
section. Define $X=2\omega/m$ and $y_B=2\omega'/m$ in the electron
rest frame. Energy conservation and the angular Jacobian give
\begin{equation}
 \frac{X}{1+X}\leq y_B\leq X,\qquad
 1-\cos\theta=2\left(\frac1{y_B}-\frac1X\right),\qquad
 \frac{\dd\cos\theta}{\dd y_B}=\frac2{y_B^2}.
\end{equation}
Thus the dimensionless squared amplitude and the cross section are
\begin{align}
 \mathcal B(X,y_B)&=\frac{B_{\mathrm{amp}}}{e_{\mathrm T}^4}
 =2\left[\frac X{y_B}+\frac{y_B}X
 -4\left(\frac1{y_B}-\frac1X\right)
 +4\left(\frac1{y_B}-\frac1X\right)^2\right],\\
 \frac{\dd\sigma_0}{\dd y_B}&=
 \frac{\pi\alpha_{\mathrm T}^2}{m^2X^2}\mathcal B(X,y_B).
\end{align}
A primitive of $\mathcal B/2$ is
\begin{equation}
 P_X(y)=\frac{y^2}{2X}
 +\left(\frac4X+\frac4{X^2}\right)y
 +\left(X-4-\frac8X\right)\ln y-\frac4y.
\end{equation}
Substitution of the two physical endpoints gives
\begin{equation}
 \sigma_0=\frac{\pi\alpha_{\mathrm T}^2}{m^2}B(X),\qquad
 B(X)=\frac2{X^2}\left[P_X(X)-P_X\!\left(\frac X{1+X}\right)\right].
\end{equation}
Expanding this endpoint difference yields the Born function stated in
Chapter 1. Its normalization has therefore been obtained from the
amplitude and phase space.

\section{Real emission and phase space}
Let $k_1,k_2$ be outgoing photon momenta. Put
$\kappa=(k,-k_1,-k_2)$ and
$E=(\slashed\epsilon,\slashed\epsilon_1^*,\slashed\epsilon_2^*)$.
The six electron-line orders give
\begin{equation}
 \mathcal M_{\gamma\gamma}=e_{\mathrm T}^3\bar u(p')\mathcal K u(p),
 \qquad
 \mathcal K=\sum_{(a,b,c)\in\mathfrak S_3}
 E_c S_0(p+\kappa_a+\kappa_b)E_b S_0(p+\kappa_a)E_a.
 \label{eq:comp-supp-double-amplitude}
\end{equation}
Here $\mathfrak S_3$ permutes $(0,1,2)$. To evaluate its unpolarized
square, sum over the two transverse polarizations of each photon,
\begin{equation}
 \overline{|\mathcal M_{\gamma\gamma}|^2}
 =\frac{e_{\mathrm T}^6}{4}\sum_{\epsilon,\epsilon_1,\epsilon_2}
 \operatorname{tr}[(\slashed p'+m)\mathcal K
                             (\slashed p+m)\overline{\mathcal K}].
\end{equation}
The factor $1/4$ averages the initial electron and photon states.
This trace includes all 36 interferences, rather than six separate
squares. Replacing any one polarization by its momentum telescopes
the six permutations to zero. Relabelling $1\leftrightarrow2$ simply
permutes the sum.

For $P=p+k$, introduce $Q=k_1+k_2$ and the standard invariant
phase space
\begin{align}
 \dd\Phi_n(P)&=(2\pi)^4\delta^{(4)}(P-\textstyle\sum_j p_j)
             \prod_j\frac{\dd^3\bm p_j}{(2\pi)^3 2p_j^0},\\
 \dd\Phi_3(P;p',k_1,k_2)&=\frac{\dd Q^2}{2\pi}
       \dd\Phi_2(P;p',Q)\dd\Phi_2(Q;k_1,k_2),\\
 \dd\Phi_2(P;a,b)&=\frac{\sqrt{\Lambda(P^2,a^2,b^2)}}
                          {32\pi^2P^2}\dd\Omega,
 \qquad\Lambda(x,y,z)=x^2+y^2+z^2-2xy-2xz-2yz.
\end{align}
The solid angle is in the rest frame of its decaying momentum.
For two massless photons this gives
\begin{equation}
 \dd\Phi_3=\frac{\sqrt{\Lambda(s,m^2,Q^2)}}{2048\pi^5s}
             \dd Q^2\dd\Omega_Q\dd\Omega_*,\qquad
 0\leq Q^2\leq(\sqrt s-m)^2.
\end{equation}
Integrating both labelled photons requires the factor $1/2!$ in
the cross section. Restricting to one photon ordering instead would
remove that divisor.

\section{The soft subtraction in the electron rest frame}
When one outgoing photon is soft with momentum $\ell$, only the
external-line denominators generate an inverse power of its energy.
For example, before mass regularization,
\begin{equation}
 \overline u(p')\gamma^\rho
 \frac{\slashed p'+\slashed\ell+m}{(p'+\ell)^2-m^2}
 =\frac{p'^\rho}{p'\cdot\ell}\overline u(p')+O((\ell^0)^0).
\end{equation}
The on-shell Dirac equation and
$\gamma^\rho\slashed p'=2p'^\rho-\slashed p'\gamma^\rho$
give the numerator $2p'^\rho$ at leading order. Emission from the
incoming leg similarly gives $-p^\rho/(p\cdot\ell)$ because its
denominator is $-2p\cdot\ell$. The other insertions are finite in
this limit. Their sum therefore gives
\begin{equation}
 \mathcal M_{\gamma\gamma}\longrightarrow
 e_{\mathrm T}\epsilon^*(\ell)\cdot J(\ell)\mathcal M_0,
 \qquad J=\frac{p'}{p'\cdot\ell}-\frac{p}{p\cdot\ell},
 \qquad \ell\cdot J=0.
\end{equation}
A common sign does not affect the rate. Define the relative soft
correction by
\begin{equation}
 \delta_{\mathrm s}(\varepsilon,\lambda)
 =e_{\mathrm T}^2\int_{E_\ell<\varepsilon}
 \frac{\dd^3\bm\ell}{(2\pi)^3\,2E_\ell}[-J(\ell)^2],
 \qquad E_\ell=\sqrt{|\bm\ell|^2+\lambda^2}.
\end{equation}
Here $\varepsilon$ is an auxiliary photon-energy boundary satisfying
$\lambda\ll\varepsilon$ and small compared with the hard energies.
Current conservation $\ell\cdot J=0$ removes longitudinal terms
from the photon polarization sum. The expression is exact in
$\lambda$ within the leading soft current. Terms beyond that current
vanish when the soft-energy boundary is removed.
In the rest frame of $p$,
choose the polar axis along $\bm p'$. Set
$E'=p'^0$, $v=|\bm p'|/E'$, $q=|\bm\ell|$,
$E_\ell=\sqrt{q^2+\lambda^2}$, $\beta=q/E_\ell$, and
$c=\widehat{\bm p'}\cdot\widehat{\bm\ell}$. Then
\begin{equation}
 -J^2=\frac1{E_\ell^2}
 \left[\frac2{1-v\beta c}-1
                  -\frac{1-v^2}{(1-v\beta c)^2}\right].
\end{equation}
The azimuthal integral is $2\pi$. For $a=v\beta$ the two polar
integrals are
\begin{equation}
 \int_{-1}^1\frac{\dd c}{1-ac}
 =\frac{2\operatorname{artanh}a}{a},\qquad
 \int_{-1}^1\frac{\dd c}{(1-ac)^2}=\frac{2}{1-a^2}.
\end{equation}
Their values at $a=0$ are understood by continuity. Substituting them
and $e_{\mathrm T}^2=4\pi\alpha_{\mathrm T}$ gives
\begin{equation}
 \begin{aligned}
 \delta_{\mathrm s}(\varepsilon,\lambda)
 =\frac{\alpha_{\mathrm T}}{2\pi}
 \int_0^{\sqrt{\varepsilon^2-\lambda^2}}
 \frac{q^2\dd q}{E_\ell^3}
 \left[\frac4{v\beta}\operatorname{artanh}(v\beta)
       -2-\frac{2(1-v^2)}{1-v^2\beta^2}\right].
 \end{aligned}
 \label{eq:comp-supp-soft-exact}
\end{equation}
This one-dimensional integral supplies the finite terms as well as
the logarithm. For $q\gg\lambda$ its bracket approaches $4(b-1)$,
where $b=\operatorname{artanh}(v)/v$ and
$v=\sqrt{1-m^4/(p\cdot p')^2}$. Therefore
\begin{equation}
 \delta_{\mathrm s}=
 \frac{2\alpha_{\mathrm T}}\pi(b-1)\ln\frac\varepsilon\lambda
 +\text{finite terms}.
\end{equation}
The soft region of the virtual graph has the same two eikonal
lines. Its interference has the opposite coefficient of
$\ln\lambda$,
\begin{equation}
 V=B_{\mathrm{amp}}\frac{2\alpha_{\mathrm T}}\pi(b-1)\ln\frac\lambda m
                   +V_{\mathrm{finite}}.
\end{equation}
Consequently $V+B_{\mathrm{amp}}\delta_{\mathrm s}$ is independent of $\lambda$
in its zero-mass limit. The finite remainder can be evaluated with
the parameter integrals above and \eqref{eq:comp-supp-soft-exact}.

For dimensionless measurement functions $F_2,F_3$, evaluated on
the two-body and three-body events respectively, require symmetry
under exchange of the two final photons
and $F_3\to F_2$ as either photon becomes soft. The observable is
\begin{equation}
 \begin{aligned}
 \sigma^{\mathrm{NLO}}[F]=\frac1{4p\cdot k}
 \lim_{\varepsilon\to0}\Bigg\{&\int\dd\Phi_2F_2
 [B_{\mathrm{amp}}+\lim_{\lambda\to0}(V+B_{\mathrm{amp}}\delta_{\mathrm s})]\\
 &+\frac12\int_{\min(k_1^0,k_2^0)>\varepsilon}
 \dd\Phi_3 F_3\overline{|\mathcal M_{\gamma\gamma}|^2}\Bigg\}
 +\frac1{m^2}\Order(\alpha_{\mathrm T}^4).
 \end{aligned}
\end{equation}
The total considered in this document has $F_2=F_3=1$.
All separator energies in this expression are in the same electron
rest frame. The two soft corners cancel the divisor $2!$ in the
soft part. The hard integral cancels the separator logarithm.
At a fixed nonzero separator, there are power corrections from
replacing the soft event by its two-body limit. Both limits are
needed before interpreting the result as an inclusive observable.
